% !TeX encoding = UTF-8
\documentclass[variant=guia,cover=institutional,mode=screen]{ua-engineering-v6-article}
% !TeX encoding = UTF-8
\UASetUniversity{Universidad Austral}
\UASetFaculty{Facultad de Ingeniería}
\UASetDepartment{Departamento de Computación}
\UASetCampus{Pilar}
\UASetCareer{Ingeniería Informática}
\UASetCommission{Comisión A}
\UASetProfessor{Equipo docente}
\UASetSemester{Segundo cuatrimestre 2026}
\UASetCourse{Sistemas Distribuidos}
\UASetDate{2026}


\UASetClassNumber{11}
\UASetClassTitle{Resiliencia distribuida}
\UASetDocKind{Guía de ejercicios}

\title{Clase 11 --- Guía de ejercicios}
\author{\UAProfessor}
\date{\UADate}

\begin{document}

\section{Indicaciones generales}
Resuelva cada ejercicio declarando supuestos, modelo de fallas, garantía y trade-off. Cuando corresponda, construya una ejecución verificable y vincule la respuesta con evidencia observable.

\section{Nivel 1: fundamentos y reconocimiento}
\begin{uaexercise}
Explique por qué una dependencia lenta puede ser más peligrosa que una dependencia completamente caída.
\end{uaexercise}

\begin{uaexercise}
Construya una cadena causal donde una degradación local termine generando una caída visible en el borde del sistema.
\end{uaexercise}

\begin{uaexercise}
Explique por qué la resiliencia distribuida no consiste en “hacer más retries”.
\end{uaexercise}

\begin{uaexercise}
Explique la diferencia entre tolerar una falla y amplificarla.
\end{uaexercise}

\begin{uaexercise}
Explique en qué casos retry es razonable y en qué casos es peligroso.
\end{uaexercise}


\section{Nivel 2: ejecuciones y fallas}
\begin{uaexercise}
Diseñe una política de retry para una operación idempotente.
\end{uaexercise}

\begin{uaexercise}
Explique por qué el exponential backoff es mejor que reintentos inmediatos constantes.
\end{uaexercise}

\begin{uaexercise}
Explique el rol del jitter en un sistema con miles de clientes concurrentes.
\end{uaexercise}

\begin{uaexercise}
Construya un ejemplo donde retries en múltiples capas provoquen una tormenta de requests.
\end{uaexercise}

\begin{uaexercise}
Defina circuit breaker y describa sus estados.
\end{uaexercise}


\section{Nivel 3: diseño y comparación}
\begin{uaexercise}
Explique con un ejemplo qué gana el sistema al abrir un breaker.
\end{uaexercise}

\begin{uaexercise}
Construya un caso donde un threshold mal calibrado de breaker empeore el comportamiento del sistema.
\end{uaexercise}

\begin{uaexercise}
Explique por qué el estado half-open es importante.
\end{uaexercise}

\begin{uaexercise}
Explique qué problema resuelve backpressure.
\end{uaexercise}

\begin{uaexercise}
Diseñe un ejemplo donde la ausencia de backpressure haga crecer una cola hasta volver inútil el sistema.
\end{uaexercise}


\section{Nivel 4: integración y defensa}
\begin{uaexercise}
Explique qué es bulkhead y cómo se implementa en términos de recursos.
\end{uaexercise}

\begin{uaexercise}
Construya un caso donde un pool compartido arrastre rutas críticas que no tenían problema.
\end{uaexercise}

\begin{uaexercise}
Defina graceful degradation y compárela con caída total.
\end{uaexercise}

\begin{uaexercise}
Diseñe tres ejemplos concretos de degradación elegante en una plataforma de e-commerce.
\end{uaexercise}

\begin{uaexercise}
Explique qué riesgos aparecen si la degradación no es explícita y semánticamente clara.
\end{uaexercise}


\end{document}
